% Section 6: Discussion
% Note: \section{} header is in main.tex

\subsection{Format vs.\ Semantic Collapse: Implications for Practitioners}

Existing diversity-preserving methods target semantic diversity as their primary objective. DIVE \cite{placeholder_dive} penalizes embedding similarity during candidate selection, DARLING \cite{placeholder_darling} co-optimizes a diversity reward alongside task quality, and GEM \cite{placeholder_gem} trains multiple policies to cover distinct output modes. All three assume that reward optimization narrows the semantic distribution.

Our results challenge this assumption. Across seven configurations and three generations, Vendi Score remains within $\pm$1\% of its initial value. The collapse that does occur is structural: opening entropy drops 1.13 bits under curation, and the near-duplicate start rate rises from 50\% to 84\%. A practitioner who monitors only Vendi Score or Self-BLEU will conclude that diversity is intact, while end users encounter responses that begin identically and follow the same template.

Users perceive repetition through surface patterns, not through embedding geometry. A system whose outputs are semantically varied but structurally uniform will feel monotonous in deployment. We recommend that diversity audits include at least one format-level metric alongside embedding-based measures. Opening entropy and near-duplicate start rate are lightweight, require no GPU, and capture a dimension of homogenization that current pipelines miss.

\subsection{Why Semantic Collapse Does Not Occur}

Section~\ref{sec:analysis} identifies three mechanisms that prevent semantic narrowing: prompt diversity anchoring, reward saturation, and generation stochasticity. Here we situate these mechanisms within the broader model collapse literature.

The theoretical conditions for model collapse \cite{placeholder_shumailov} assume that a model trains on its own unfiltered outputs, with each generation's full distribution as the next generation's training data. Under this regime, tail modes lose probability mass across iterations and the distribution contracts. Reward-curated pipelines violate two of these preconditions. First, the training distribution is a filtered subset selected by an external reward signal, not the model's full output. The reward model introduces structure independent of the generator's biases. Second, the prompt set provides an external scaffold: 2,200 diverse instructions force broad semantic coverage at every generation, regardless of how previous outputs were filtered.

Reward saturation further limits compounding. The score gap between the top-1 candidate and the pool average is approximately 0.04 on a normalized scale, so curation selects responses only marginally preferred by the reward model rather than responses from a single dominant mode.

These observations do not imply that semantic collapse can never occur under reward curation. A narrow prompt domain (e.g., single-task coding) would remove the diversity anchor, a highly discriminative reward model with large score gaps would strengthen the selection filter, and deterministic decoding ($T{=}0$) would eliminate generation stochasticity. Our claim is bounded: under the conditions tested, the theoretical preconditions for semantic collapse are not met.

\subsection{LoRA Evaluation Caveat}

The direction reversal we observe between LoRA and Full-FT configurations (Section~\ref{sec:results}) has implications beyond this study. LoRA has become the default adaptation method for alignment research due to its low computational cost \cite{placeholder_hu2022lora}. Our results show that LoRA's rank-16 bottleneck (1.67\% of parameters) imposes an upper bound on how much distributional structure the model can absorb from curated data. Under LoRA, curation suppresses Vendi Score by 3.40 percentage points relative to the random baseline. Under Full-FT at the same learning rate, this gap shrinks to 0.46 points. At a higher learning rate, the effect reverses entirely: curation \emph{promotes} diversity by 0.63 points.

This pattern suggests that conclusions drawn from LoRA-based evaluations of curation or reward optimization may conflate capacity constraints with method properties. We do not argue that LoRA is unsuitable for alignment work. Rather, when the research question concerns distributional effects of training data selection, evaluation under full fine-tuning is necessary to separate capacity artifacts from genuine curation dynamics. Studies that report diversity suppression under LoRA alone risk overstating the cost of reward-based selection.

\subsection{Reward Model Alignment as a Limiting Factor}

A supplementary analysis reveals that ArmoRM's preferences correlate with held-out judge scores at $r = 0.63$--$0.77$, depending on the evaluation subset. This partial alignment has two consequences.

First, it explains the near-zero quality improvement across generations: curated training data carries a reward signal only partially informative of true quality, so the model learns to satisfy ArmoRM's biases rather than to produce genuinely better outputs. Second, partial alignment may also explain why curation does not collapse semantic diversity. A perfectly aligned reward model with sharp preferences could drive the generator toward a narrow set of optimal responses; a noisy signal, by contrast, selects a broader band of outputs that happen to score well, preserving distributional breadth through imprecision.

If the same reward model that curates training data also serves as the quality metric, improvements are circular. Independent evaluation through held-out judges, human ratings, or orthogonal automatic metrics remains essential for any iterative self-training pipeline.
